\documentclass[11pt,a4paper]{article}
\usepackage[utf8]{inputenc}
\usepackage{amsmath,amssymb,amsthm}
\usepackage[margin=1in]{geometry}
\usepackage{enumitem}
\usepackage{hyperref}

\newtheorem{theorem}{Theorem}
\newtheorem{proposition}[theorem]{Proposition}
\newtheorem{definition}[theorem]{Definition}

\title{arXiv Report: Existence, Properties, and Parametric Inference for\\Possibly Hyperuniform Gibbs Perturbed Lattices}
\author{Auto-generated report for arXiv:2603.01858\\J.-F.\ Coeurjolly and C.\ Renaud-Chan}
\date{March 4, 2026}

\begin{document}
\maketitle

%% ----------------------------------------------------------------
\section{Core Question}

Can one build a flexible class of point processes that (i) retains a lattice backbone, (ii) allows Gibbs-type interactions (including hard-core repulsion), and (iii) may be \emph{hyperuniform} --- a property that neither classical Gibbs point processes nor simple perturbed lattices can simultaneously achieve?

%% ----------------------------------------------------------------
\section{Exact Main Result}

\begin{theorem}[Existence, Theorem~1 of the paper]
Let $\mathcal{L}\subset\mathbb{R}^d$ be a full-rank lattice with fundamental domain $D$.
Assume the Hamiltonian $H$ satisfies \textup{(H1)--(H4)} (stability, non-degeneracy, heredity, translation-invariance) and either
\begin{itemize}[nosep]
  \item \textup{(H5.1)} finite range $+$ \textup{(Q1)[1]} (exponential moment of order $1$ for the move distribution $Q$), or
  \item \textup{(H5.2)} pairwise summability $+$ \textup{(Q2)} (bounded support of $Q$).
\end{itemize}
Then the set $\mathcal{G}_{\mathcal{L}}(H,Q)$ of Gibbs measures is non-empty, and the associated \emph{Gibbs perturbed lattice}
\[
  \Gamma = \{i + X_i + U : i\in\mathcal{L}\}, \qquad X\sim P\in\mathcal{G}_{\mathcal{L}}(H,Q),\quad U\sim\mathrm{Unif}(D),
\]
is a stationary point process on $\mathbb{R}^d$.  Moreover, for $d\le 2$, under strengthened moment/boundedness conditions on $Q$, $\Gamma$ is \textbf{hyperuniform}:
\[
  \frac{\mathrm{Var}[\#\Gamma\cap B(0,R)]}{|B(0,R)|}\;\longrightarrow\;0 \quad\text{as }R\to\infty.
\]
\end{theorem}

%% ----------------------------------------------------------------
\section{Concrete Example}

\paragraph{Strauss hard-core perturbed lattice.}
Take $\mathcal{L}=\mathbb{Z}^2$ and define the pairwise Hamiltonian
\[
  H(\gamma)=\sum_{\{x,y\}\subset\gamma}\Bigl[A\,\mathbf{1}_{[r,R]}(|x-y|) + \infty\cdot\mathbf{1}_{[0,r)}(|x-y|)\Bigr],
\]
with interaction range $R>0$, hard-core distance $r\ge 0$, and energy parameter $A\ge 0$.
Choose the move distribution $Q$ as the uniform distribution on a ball $B(0,\varepsilon)$ with $\varepsilon<(\delta-2r)/2$ (where $\delta$ is the lattice spacing).

\smallskip\noindent
The resulting Gibbs perturbed lattice:
\begin{itemize}[nosep]
  \item has a \emph{guaranteed} hard-core: no two points can be closer than $r$;
  \item exhibits Strauss-type soft repulsion at distances in $[r,R]$;
  \item is \textbf{hyperuniform} in $d=2$ (by Proposition~3.6 of the paper).
\end{itemize}
None of these three properties can be simultaneously achieved by a standard Gibbs point process on $\mathbb{R}^2$ or a simple i.i.d.\ perturbed lattice.

%% ----------------------------------------------------------------
\section{Intuition and Setup}

\paragraph{Key objects.}
\begin{itemize}[nosep]
  \item \textbf{Lattice} $\mathcal{L}\subset\mathbb{R}^d$: a full-rank discrete subgroup generated by $d$ linearly independent vectors $a_1,\dots,a_d$.  Fundamental domain $D=\bigl\{\sum t_i a_i : t_i\in[0,1)\bigr\}$.  Lattice spacing $\delta=\min\{|i|:i\in\mathcal{L}\setminus\{0\}\}$.
  \item \textbf{Move distribution} $Q$: a probability measure on $\mathbb{R}^d$ controlling how far each particle can be displaced from its lattice site.
  \item \textbf{Perturbed lattice}: $\Gamma=\{i+X_i+U:i\in\mathcal{L}\}$ where each $X_i\sim Q$ (not necessarily i.i.d.---the joint law is governed by $P$) and $U\sim\mathrm{Unif}(D)$ is an independent global shift that stationarizes $\Gamma$.
  \item \textbf{Hamiltonian} $H:\mathcal{C}_f\to\mathbb{R}\cup\{\infty\}$: measures interaction energy of a finite point configuration.  The local energy of a point $x$ given the rest of the configuration $\gamma$ is $h(x,\gamma)=\lim_{n\to\infty}[H(\{x\}\cup\gamma_{[-n,n]^d})-H(\gamma_{[-n,n]^d})]$.
  \item \textbf{Gibbs measure on $\mathcal{L}$}: the joint law $P$ of $(X_i)_{i\in\mathcal{L}}$ satisfying DLR equations~\eqref{eq:dlr} below.
\end{itemize}

\paragraph{Why a new model class?}
Known hyperuniform processes (Ginibre, determinantal, i.i.d.\ perturbed lattices) lack flexibility: one cannot impose hard-core constraints.
Conversely, Coeurjolly--Flint--Laverny (2024) showed that most Gibbs point processes on $\mathbb{R}^d$ are \emph{not} hyperuniform.
The authors bridge this gap by making the Hamiltonian act on the \emph{perturbed locations} $\{i+X_i\}$ rather than on the marks $\{X_i\}$ alone, creating a model that is simultaneously:
\begin{enumerate}[nosep]
  \item lattice-based (hence potentially hyperuniform),
  \item Gibbs-type (hence flexible, with tunable interactions).
\end{enumerate}

%% ----------------------------------------------------------------
\section{Main Results (Summary)}

\begin{itemize}
  \item \textbf{Existence (Theorem~1).}  Under stability, non-degeneracy, heredity, translation-invariance of $H$, plus a range/moment condition, infinite-volume Gibbs measures on $\mathcal{L}$ exist and yield stationary point processes.

  \item \textbf{Hyperuniformity (Proposition~3.6).}
    \begin{itemize}[nosep]
      \item[(i)] Under (H5.1)+(Q1)[$d$] or (H5.2)+(Q2): \emph{some} Gibbs perturbed lattice in $d\le 2$ is hyperuniform.
      \item[(ii)] Under (H5.3)+(Q3) (bounded local energy + $d$-th moment): \emph{every} Gibbs perturbed lattice in $d\le 2$ is hyperuniform.
    \end{itemize}

  \item \textbf{DLR-type equations (Propositions~3.7--3.9).}  Three families of equilibrium equations (first-order, second-order, variational) analogous to GNZ equations for classical Gibbs point processes.

  \item \textbf{Parametric inference (Theorem~2).}
    For exponential-family models observed in an expanding window $W_n$:
    \begin{itemize}[nosep]
      \item A Takacs--Fiksel estimator $\hat\theta_n$ is defined, encompassing pseudo-likelihood and variational estimators as special cases.
      \item $\hat\theta_n\to\theta^\star$ a.s.\ (consistency).
      \item $\sqrt{|W_n|}(\hat\theta_n-\theta^\star)\Rightarrow\mathcal{N}(0,\Sigma)$ (asymptotic normality), where $\Sigma$ is an explicit sandwich-type covariance matrix.
    \end{itemize}

  \item \textbf{Two observation frameworks:}
    (1) both lattice sites $i$ and displacements $X_i$ are observed;
    (2) only the perturbed locations $i+X_i$ are observed (harder; handled via the variational equation for specific move distributions).
\end{itemize}

%% ----------------------------------------------------------------
\section{Proof Sketch}

\subsection*{Step 1: Finite-volume Gibbs measures}
On a finite sub-lattice $\Lambda\subset\mathcal{L}$, define the ``lifted'' Hamiltonian $\tilde{H}(x_\Lambda)=H(\{i+x_i:i\in\Lambda\})$.
The Gibbs measure on $\Lambda$ is
\[
  P_\Lambda(dx_\Lambda) = \frac{1}{Z_\Lambda}\,e^{-\tilde{H}(x_\Lambda)}\,Q^{\otimes\Lambda}(dx_\Lambda).
\]
Stability (H1) and non-degeneracy (H2) guarantee $Z_\Lambda\in(0,\infty)$.
DLR equations on the finite lattice follow by elementary conditioning (Proposition~3.3).

\subsection*{Step 2: Infinite-volume limit (specific entropy)}
The authors adapt the \emph{specific entropy} approach of Georgii (2011).
For an increasing sequence $\Lambda_n\uparrow\mathcal{L}$, they show:
\begin{enumerate}[nosep]
  \item The specific entropy $h(P\,|\,Q^{\otimes\mathcal{L}})$ is well-defined and lower-semicontinuous.
  \item The specific energy $e(P)=\lim_{n}\frac{1}{|\Lambda_n|}E_P[\tilde{H}(X_{\Lambda_n})]$ exists under the stability/range assumptions.
  \item The free-energy functional $f(P)=h(P\,|\,Q^{\otimes\mathcal{L}})+e(P)$ is minimized by some $P^\star$, which is then shown to satisfy the DLR equations~\eqref{eq:dlr}.
\end{enumerate}

The DLR equations for the infinite lattice are:
\begin{equation}\label{eq:dlr}
  P(dx'_\Delta\mid x_{\Delta^c}) = \frac{1}{Z_\Delta(x_{\Delta^c})}\,e^{-\tilde{H}_\Delta(x'_\Delta\cup x_{\Delta^c})}\,Q^{\otimes\Delta}(dx'_\Delta), \quad \forall\;\text{finite }\Delta\subset\mathcal{L}.
\end{equation}

\subsection*{Step 3: Hyperuniformity}
The key input is the result of Klatt--Last--Yogeshwaran (2020): a perturbed lattice $\{i+X_i+U\}$ is hyperuniform in $d\le 2$ whenever $E[|X_0|^d]<\infty$.
\begin{itemize}[nosep]
  \item Under (H5.1)+(Q1)[$d$]: the specific-entropy minimizer $P^\star$ inherits finite $d$-th moment from $Q$, giving at least one hyperuniform Gibbs perturbed lattice.
  \item Under (H5.3)+(Q3): the bounded local energy allows a uniform integrability argument showing $E_P[|X_0|^d]<\infty$ for \emph{every} $P\in\mathcal{G}_\mathcal{L}(H,Q)$.
\end{itemize}

\subsection*{Step 4: DLR-type (GNZ-type) equations}
Starting from the single-site DLR equation and applying it with test functions $f$, the authors derive:
\begin{itemize}[nosep]
  \item \emph{First-order} equation (Proposition~3.7): an expectation identity involving the normalized Papangelou intensity $\Lambda(i,x,\gamma)=\lambda(i,x,\gamma)/\int_Q\lambda(i,y,\gamma)\,dy$.
  \item \emph{Second-order} equation (Proposition~3.8): obtained by iterating the single-site DLR twice.
  \item \emph{Variational} equation (Proposition~3.9): integration by parts on the mark space using the DLR kernel, yielding an identity involving $\nabla h$ and $\nabla q$.
\end{itemize}

\subsection*{Step 5: Parametric inference}
For exponential-family models $-\log\lambda(i,x,\gamma;\theta)=\theta^\top S(i,x,\gamma)$:
\begin{enumerate}[nosep]
  \item Define the Takacs--Fiksel contrast $\mathrm{TF}_n(f,\Gamma;\theta)=\sum_{l=1}^L[\mathrm{DLR}_n(f_l,\Gamma;\theta)]^2$ where $\mathrm{DLR}_n$ is the empirical DLR residual with border correction.
  \item \textbf{Consistency:} Use a novel ergodic theorem for the empirical DLR functionals (adapting standard lattice ergodic theory to the continuous-mark, perturbed-location setting) to show $\mathrm{TF}_n/|W_n|\to\mathrm{TF}(\theta)$ uniformly, which has a unique minimum at $\theta^\star$.
  \item \textbf{Asymptotic normality:} Taylor-expand $\mathrm{TF}_n^{(1)}$ around $\theta^\star$; use a CLT for the empirical DLR functionals (proved via cumulant bounds exploiting finite range) to obtain $\sqrt{|W_n|}(\hat\theta_n-\theta^\star)\Rightarrow\mathcal{N}(0,\Sigma)$ with sandwich covariance $\Sigma=I(\theta^\star)^{-1}V(\theta^\star)I(\theta^\star)^{-1}$.
\end{enumerate}

%% ----------------------------------------------------------------
\section*{Reference}
J.-F.\ Coeurjolly and C.\ Renaud-Chan, \emph{Existence, properties, and parametric inference for possibly hyperuniform Gibbs perturbed lattices}, arXiv:2603.01858, March 2026.

\end{document}
